\documentclass[10pt,letterpaper]{article}
\usepackage[margin=0.78in]{geometry}
\usepackage{amsmath,amssymb,booktabs,graphicx,xcolor,microtype,enumitem}
\usepackage[colorlinks=true,urlcolor=blue!55!black,linkcolor=blue!55!black,citecolor=blue!55!black]{hyperref}
\usepackage{xurl}
\definecolor{navy}{HTML}{203C56}
\setlength{\parindent}{0pt}\setlength{\parskip}{5pt}
\setlist{nosep,leftmargin=*}
\begin{document}
{\Large\bfseries\color{navy} Cosmic muons through the ESA test stand}\par
{\large A rate estimate with and without 3 ft of concrete}\par
18 September 2026 \hfill CAD-registered geometry: \texttt{eae30ac8f}
\vspace{4pt}\hrule\vspace{7pt}

\textbf{Result.} The current geometry gives approximately \textbf{1.41 muons/hour without an overhead slab} and \textbf{1.20 muons/hour with 3 ft of concrete}. Both numbers require a single muon to intersect all 26 active layers. They are conditional predictions from a sea-level spectrum and a straight-line continuous-slowing-down model, before detector efficiency; they are not measured rates or a full Geant4 transport result.

\begin{center}
\begin{tabular}{lrrrr}
\toprule
Overhead shielding & Muons/hour & Muons/day & Mean interval & $K_{\min}$ at surface\\
\midrule
None & 1.41 & 33.8 & 42.6 min & 0.104 GeV\\
3 ft concrete & 1.20 & 28.8 & 50.1 min & 0.470 GeV\\
\bottomrule
\end{tabular}
\end{center}
Here $K_{\min}$ is the acceptance-weighted mean kinetic energy needed to emerge from the modelled detector with at least 10 MeV remaining. That small residual is a range-table reliability floor, not an electronics threshold. The concrete reduces this model's rate by approximately \textbf{15\%}.

\section*{What counts as going through everything?}
A trajectory must have nonzero path length in an active volume in \emph{every layer} listed below. Within a segmented layer, hitting any one bar is sufficient. A track need not hit all bars in that layer.
\begin{center}
\begin{tabular}{lrrl}
\toprule
Subsystem & Required layers & Active volumes & Geometry used\\
\midrule
Top HCal & 2 & 24 & Two crossed layers, 12 bars each\\
Trigger counters & 6 & 72 & Three stations, two staggered layers each\\
Silicon tracker & 4 & 4 & Active daughters; inactive borders excluded\\
LYSO target & 2 & 33 & Two staggered bar arrays\\
HCal above ECal & 4 & 32 & Four crossed layers, 8 bars each\\
ECal silicon & 4 & 4 & Four hexagonal active silicon layers\\
HCal below ECal & 4 & 48 & Four crossed layers, 12 bars each\\
\midrule
Total & 26 & 217 & Bars and sensors are distinct volumes\\
\bottomrule
\end{tabular}
\end{center}
The STEP model supplies the placement framework; ESA \texttt{ldmx-reduced-v3} supplies the ECal, trigger, tracker and target internals. HCal dimensions follow the test-stand branch. The inferred HCal bar counts and detector registration still need hardware confirmation. The support renderings are not additional material in the GDML used here.

\section*{The geometric bottleneck}
The trigger bars are only 30 mm wide in $X$, while the HCal spans metres. After requiring the six trigger and two LYSO layers, the tracker and HCal do not remove further tracks in this integration; the displaced ECal footprint reduces the remaining rate by about 36\%. The accepted rays are close to vertical: 90\% are within 4.82$^\circ$ and 99\% within 6.28$^\circ$ under the reference angular model.

For the simple reference intensity $I(\theta)=70\cos^2\theta\ \mathrm{m}^{-2}\mathrm{s}^{-1}\mathrm{sr}^{-1}$, associated with sea-level muons above $1\ \mathrm{GeV}/c$\cite{pdg},
\[
G_2=\int dA\,d\Omega\,\cos^3\theta\,\mathbf{1}_{\mathrm{all}}=4.0021\times10^{-6}\ \mathrm{m}^2\mathrm{sr},\qquad
R_{p>1}^{\mathrm{geom}}=1.0085\pm0.0004\ \mathrm{hour}^{-1}.
\]
The quoted error is numerical integration error only. The headline rates are larger because the energy-dependent calculation also includes lower-energy muons that have enough range to cross the selected materials.

\newpage
\section*{Geometry and shielding illustrations}
\begin{center}\includegraphics[width=\linewidth]{geometry_acceptance.pdf}\end{center}
\textbf{Figure 1.} Two projections of 32 accepted sample trajectories. Coloured segments mark the actual layer centres and transverse extents; their line thickness is enlarged for legibility. HCal bars extend beyond the displayed transverse window. The ECal offset and narrow trigger/LYSO corridor restrict the accepted directions. These are straight geometric rays, not simulated scattering histories.

\vspace{7pt}
\begin{center}\includegraphics[width=\linewidth]{rate_comparison.pdf}\end{center}
\textbf{Figure 2.} The continuous-loss estimate versus concrete thickness, and the angular distribution of muons that satisfy all 26 layer requirements. The 3 ft comparison assumes a uniform horizontal slab that covers the stand and all accepted directions, with concrete density $2.30\ \mathrm{g\,cm^{-3}}$\cite{concrete}. Roof height and openings are not modelled. No efficiency correction or calibrated uncertainty band is implied.

\newpage
\section*{Calculation and checks}
\textbf{Geometry integration.} Two points are sampled on horizontal planes through the upper trigger and lower LYSO layer, separated by $D=383.4$ mm. Sampling rectangles are slightly expanded to include tracks that clip a finite-thickness anchor volume. If $\boldsymbol{u}=({\bf b}-{\bf a})/D$, then
\[
R=\int d^2a\,d^2b\;\frac{I(\theta)\cos^4\theta}{D^2}\,
\mathbf{1}_{\mathrm{all}},\qquad \cos\theta=(1+|\boldsymbol{u}|^2)^{-1/2}.
\]
This includes both the projected-area factor and the solid-angle Jacobian. For $I=I_0\cos^2\theta$, each ray carries a $\cos^6\theta$ weight. Box and hexagonal-prism intersections retain the physical bar gaps, stagger, tracker rotations and finite layer thickness. Eight independently scrambled Sobol sequences of $2^{18}$ point pairs give 2,097,152 trials. An independent pseudorandom estimate agrees within its numerical error. ROOT's original GDML solids agree with all 4,992 individual layer decisions in a separate 192-ray test.

\textbf{Spectrum.} The energy-dependent model uses Guan et al.'s modified Gaisser spectrum\cite{guan}, with total energy $E$ in GeV and $j$ in $\mathrm{cm^{-2}s^{-1}sr^{-1}GeV^{-1}}$:
\[
j(E,\theta)=C\,0.14\left[E+\frac{3.64}{(\cos\theta^*)^{1.29}}\right]^{-2.7}
\left[\frac{1}{1+1.1E\cos\theta^*/115}+\frac{0.054}{1+1.1E\cos\theta^*/850}\right].
\]
The factor $C=1.1595$ sets the vertical integral above $p=1\ \mathrm{GeV}/c$ to $70\ \mathrm{m^{-2}s^{-1}sr^{-1}}$. The published curvature correction for $\theta^*$ is implemented in the supplied script. The flux above each ray's required surface energy replaces the reference $70\cos^2\theta$ weight. The resulting lower-energy contribution is model dependent.

\textbf{Stopping model.} The supplied 3 ft thickness is $91.44$ cm, or $210.312\ \mathrm{g\,cm^{-2}}=2.10312$ m water equivalent at the assumed density. ROOT chords and the GDML densities give a mean vertical detector column of $32.70\ \mathrm{g\,cm^{-2}}$ over 96 accepted rays, of which roughly $20.6\ \mathrm{g\,cm^{-2}}$ is HCal scintillator. For this estimate, the whole detector column is represented by a polystyrene CSDA range\cite{poly}; the roof uses the concrete CSDA table\cite{concrete}. This is an approximation to the mixed detector materials, not material-by-material transport. With tabulated range functions $\mathcal R_p$ and $\mathcal R_c$,
\[
K_d=\mathcal R_p^{-1}\!\left[\mathcal R_p(10\ \mathrm{MeV})+X_d\sec\theta\right],\quad
K_s=\mathcal R_c^{-1}\!\left[\mathcal R_c(K_d)+X_c\sec\theta\right].
\]
The flux integral starts at $E=K_s+m_\mu$. With no roof, $K_s=K_d$. Doubling the detector column changes the 3 ft prediction from 1.20 to 1.17/hour; requiring 100 MeV at the final exit gives 1.17/hour.

\textbf{Interpretation.} The recorded coincidence rate will also depend on the light-yield thresholds, dead channels, trigger logic and live time. Multiple scattering, decay and stochastic energy losses are omitted. Grazing tracks can count here: requiring a hit at every layer's midplane instead reduces the reference rate from 1.0085 to 0.853/hour, showing that the signal threshold matters. The flux normalization itself is not exact; the PDG review notes measurements 10--15\% below the reference value\cite{pdg}. No combined systematic uncertainty or local flux calibration is claimed. The estimate is for cosmic muons of both charges, not all cosmic-ray species or shower coincidences.

\textbf{Reproduction.} The folder contains the geometry snapshot, scripts, numerical tables, range data, figures and this \LaTeX{} source. Run \texttt{build\_report.sh} with Python (NumPy, SciPy, Matplotlib), ROOT and Tectonic available. The \texttt{SHA256SUMS} file records artifact identities. The original STEP SHA-256 and all source GDML hashes are preserved in \texttt{geometry/provenance.json}.

\begin{thebibliography}{4}\small\setlength{\itemsep}{0pt}
\bibitem{pdg} Particle Data Group, \emph{Cosmic Rays}, 2022 review, Sec. 30.3.1.
\url{https://pdg.lbl.gov/2022/reviews/rpp2022-rev-cosmic-rays.pdf}.
\bibitem{guan} M. Guan et al., \emph{A parametrization of the cosmic-ray muon flux at sea-level}, arXiv:1509.06176 (2015), Eqs. (2)--(3). \url{https://arxiv.org/abs/1509.06176}.
\bibitem{concrete} PDG, \emph{Shielding concrete: properties and muon CSDA range}.
\url{https://pdg.lbl.gov/2022/AtomicNuclearProperties/HTML/shielding_concrete.html}.
\bibitem{poly} PDG, \emph{Muon stopping power and range in polystyrene}.
\url{https://pdg.lbl.gov/2019/AtomicNuclearProperties/MUE/muE_polystyrene.txt}.
\end{thebibliography}
\clearpage
\section*{Raising the TS/tracker/LYSO assembly by 0.20 m}
All three trigger-scintillator (TS) stations, all four tracker sensors and both
LYSO layers move together by $+200$ mm in CAD $Z$. Their transverse positions,
rotations, internal spacings and segmentation stay the same. ECal and all three
HCal groups remain fixed. A muon must still cross every one of the 26 active layers.

\textbf{Raising this assembly alone reduces the estimated rate by 10.2\%.}
The top-TS to lower-LYSO separation stays at $383.4$ mm, preserving the internal
collimation, but the distance to the fixed, offset ECal grows by 200 mm. Fewer
accepted directions also cross every ECal layer. The geometric
$p>1\ \mathrm{GeV}/c$ reference changes from 1.0085 to
$0.9059\pm0.0003$ muons/hour; the error is numerical only.

\begin{center}
\begin{tabular}{lrrr}
\toprule
Overhead shielding & Original/hour & Raised/hour & Change\\
\midrule
None & 1.407 & 1.264 & $-10.2\%$\\
3 ft concrete & 1.199 & 1.077 & $-10.2\%$\\
\bottomrule
\end{tabular}
\end{center}

\subsection*{Centering the ECal as well}
Keep the TS/tracker/LYSO assembly raised and translate the entire ECal in $X,Y$
so that the centre of its combined active-silicon envelope lies on the LYSO/TS
axis. This means $\Delta X=+46.56$ mm and $\Delta Y=-46.08$ mm; ECal height,
rotations and internal layer offsets stay unchanged. All HCal positions stay fixed.
\textbf{Centering increases the rate by 54.2\% relative to raising alone},
and by 38.5\% relative to the original layout.

\begin{center}
\begin{tabular}{lrrr}
\toprule
Overhead shielding & Raised + centred/hour & Per day & Gain over raised\\
\midrule
None & 1.949 & 46.8 & $+54.2\%$\\
3 ft concrete & 1.660 & 39.8 & $+54.2\%$\\
\bottomrule
\end{tabular}
\end{center}

\begin{center}
\includegraphics[width=\linewidth]{shift_up_200mm/geometry_shift_comparison.pdf}
\end{center}
\textbf{Figure 3.} The three layouts with 24 accepted straight rays each, in
$X$--$Z$. TS, tracker and LYSO move upward together in B and C; ECal is also
centred in C. Its $Y$ shift is not visible in this projection. Line thicknesses
are enlarged, and HCal extends outside the transverse window.

\textbf{Checks and limits.} Each variation uses eight Sobol scrambles of $2^{18}$
trials. ROOT agrees with all 4,992 layer decisions in each 192-ray test, including
109 translated active solids for B and 113 for C. Both independent pseudorandom
checks agree within 0.16 numerical standard errors. The spectrum and roof
model are unchanged. To isolate placement effects, the detector stopping proxy
stays at the original $32.70\ \mathrm{g\,cm^{-2}}$ vertical column; it is not
remeasured. Scattering, detector inefficiency and mechanical clearances remain
unmodelled. These are rate studies, not validated mechanical or transport geometries.


\clearpage
\section*{Centering ECal without raising the assembly}
Keep TS, tracker, LYSO and all HCal groups at their \emph{original heights}.
Move only ECal by $\Delta X=+46.56$ mm and $\Delta Y=-46.08$ mm, using the same
active-envelope alignment as on the previous page. ECal height, rotations and
internal layer offsets stay unchanged. All 26 active layers are still required.

\textbf{Centering alone gives 2.163 muons/hour without a roof and 1.842 with
3 ft of concrete.} This is a 53.7\% increase over the original layout, and
11.0\% above the raised-and-centred configuration. In this geometry,
raising the assembly is not needed to obtain the gain from ECal alignment.

\subsection*{Removing ECal}
The second test physically removes the entire ECal subtree, including its
passive materials, while keeping every other subsystem at its original position.
The coincidence now requires \textbf{22 active layers}: six TS, four tracker,
two LYSO and ten HCal. It no longer tests passage through an ECal.

\begin{center}
\begin{tabular}{lrrrr}
\toprule
Configuration & Layers & No roof/hour & Concrete/hour & No roof/day\\
\midrule
Original & 26 & 1.407 & 1.199 & 33.8\\
ECal centred only & 26 & 2.163 & 1.842 & 51.9\\
ECal removed & 22 & 2.214 & 1.883 & 53.1\\
\bottomrule
\end{tabular}
\end{center}
Here ``concrete'' means the same 3 ft uniform overhead slab. The 22-layer
geometric reference is $1.5730\pm0.0002$ muons/hour for $p>1$ GeV/$c$;
centering ECal retains 98.5\% of that geometrical coincidence rate while
still requiring all four ECal layers.

\begin{center}\includegraphics[width=\linewidth]{ecal_comparison/all_layout_rates.pdf}\end{center}
\textbf{Figure 4.} All five configurations under the same incident-spectrum
model. The last row uses the relaxed 22-layer requirement and excludes ECal
material. The raised rows move TS, tracker and LYSO upward by 200 mm;
the two new cases use their original positions.

\textbf{Material treatment.} The centred-only case retains the original
$32.70\ \mathrm{g\,cm^{-2}}$ stopping proxy, consistent with the earlier placement
studies. For ECal removal, ROOT chords for 96 accepted rays give
$24.32\ \mathrm{g\,cm^{-2}}$ for the remaining detector. The same polystyrene
CSDA approximation represents that reduced column. At the original column,
the 22-layer angular acceptance alone would give 2.195 and 1.869
muons/hour; removing ECal material raises them to the values in the table.

\textbf{Checks and limits.} Eight Sobol scrambles of $2^{18}$ trials were run
for each new case, with independent pseudorandom checks. ROOT agrees with all
4,992 layer decisions for centering and all 4,224 for removal in 192-ray tests.
Scattering, straggling, decay, detector inefficiency and mechanical-fit validation
remain outside this estimate. The scripts and results are in
\texttt{ecal\_comparison/}; \texttt{build\_report.sh} rebuilds these sections too.

\end{document}
